\documentclass[a4paper,12pt]{article}

\usepackage[T2A]{fontenc}
\usepackage[utf8]{inputenc}
\usepackage[russian]{babel}
\usepackage{amsmath,amssymb,mathtools}
\usepackage{helvet}
\renewcommand{\familydefault}{\sfdefault}
\usepackage{icomma}
\usepackage{geometry}
\usepackage{xcolor}
\usepackage{tikz}
\usetikzlibrary{arrows.meta,positioning,shapes.geometric}
\usepackage{tabularx,booktabs}
\usepackage{enumitem}
\usepackage{hyperref}
\usepackage{parskip}
\usepackage{fancyhdr}
\usepackage{setspace}

% ============================================================
% Проверено реальной компиляцией pdfLaTeX:
% ошибок: 0
% предупреждений: 0
% ============================================================

\geometry{
  left=20mm,
  right=20mm,
  top=20mm,
  bottom=22mm
}

\onehalfspacing
\setlength{\parindent}{0pt}
\setlength{\headheight}{25pt}
\setlength{\emergencystretch}{2em}
\renewcommand{\arraystretch}{1.35}

% ============================================================
% Палитра
% ============================================================

\definecolor{accent}{HTML}{008080}
\definecolor{darkaccent}{HTML}{004D4D}
\definecolor{textgray}{HTML}{333333}

\hypersetup{
  colorlinks=true,
  linkcolor=darkaccent,
  urlcolor=darkaccent
}

% ============================================================
% Колонтитулы
% ============================================================

\pagestyle{fancy}
\fancyhf{}

\fancyhead[L]{\small\color{textgray}Математика, 7 класс}
\fancyhead[R]{\small\color{textgray}Кирилл Зиновьев}
\fancyfoot[C]{\small\color{textgray}\thepage}

\renewcommand{\headrulewidth}{0.4pt}
\renewcommand{\headrule}{%
  \hbox to\headwidth{%
    \color{accent}%
    \leaders\hrule height \headrulewidth\hfill
  }%
}

% ============================================================
% Списки
% ============================================================

\setlist[itemize]{
  leftmargin=1.4em,
  itemsep=0.25em,
  topsep=0.35em
}

\setlist[enumerate]{
  leftmargin=1.7em,
  itemsep=0.45em,
  topsep=0.4em
}

% ============================================================
% Вспомогательные команды
% ============================================================

\newcommand{\topicline}{%
  \par
  \noindent
  \color{accent}%
  \rule{\linewidth}{0.6pt}%
  \color{black}%
  \par
}

\newcommand{\term}[1]{%
  \textbf{\color{darkaccent}#1}%
}

\begin{document}

% ============================================================
% ТИТУЛЬНЫЙ ЛИСТ
% ============================================================

\begin{titlepage}
\thispagestyle{empty}

\begin{center}

\vspace*{1.8cm}

{\large\color{textgray}
МАТЕМАТИКА · 7 КЛАСС
\par}

\vspace{1.2cm}

{\Huge\bfseries\color{darkaccent}
Чек-лист
\par}

\vspace{0.45cm}

{\LARGE\bfseries
Раскрытие скобок, приведение подобных\\[0.15cm]
и подстановка значений
\par}

\vspace{1cm}

\topicline

\vspace{0.9cm}

{\large
Ключевые навыки занятия:
\par}

\vspace{0.3cm}

{\large
подобные слагаемые · знаки перед скобками ·\\
распределительный закон · упрощение выражений ·\\
контроль подстановки
\par}

\vfill

{\large
Кирилл Зиновьев
\par}

\vspace{0.35cm}

{\normalsize
\today
\par}

\vspace{1cm}

{\normalsize
Лёвин Артём Александрович эксклюзивно для Кирилла Зиновьева
\par}

\vspace{1.4cm}

\begin{tikzpicture}[x=1cm,y=1cm]
  \draw[
    accent,
    line width=1.2pt
  ]
  (-6,0)
  .. controls (-3.8,1.0) and (-1.7,-0.8) .. (0,0)
  .. controls (1.8,0.8) and (3.8,-1.0) .. (6,0);
\end{tikzpicture}

\end{center}

\end{titlepage}

% ============================================================
% ОСНОВНОЙ ЧЕК-ЛИСТ
% ============================================================

\section*{1. Подобные слагаемые}
\topicline

\term{Подобные слагаемые} --- слагаемые с одинаковой
буквенной частью, включая одинаковые степени переменных.

Например,

\[
7a \text{ и } -9a,
\qquad
4x^2y \text{ и } -3x^2y
\]

--- подобные слагаемые.

Выражения

\[
3x \text{ и } 3x^2,
\qquad
2ab \text{ и } 2a
\]

имеют разные буквенные части.

При приведении подобных складывают или вычитают
только коэффициенты:

\[
7a-9a-8b+2b
=
(7-9)a+(-8+2)b
=
-2a-6b.
\]

Числа без букв также образуют отдельную группу
подобных слагаемых:

\[
4x-7+3x+2
=
7x-5.
\]

\textbf{Проверка:}
буквенная часть после приведения подобных сохраняется.

% ============================================================

\section*{2. Как раскрывать скобки}
\topicline

Основа --- распределительный закон:

\[
k(a+b)=ka+kb,
\qquad
k(a-b)=ka-kb.
\]

Знак «минус» перед скобками означает умножение
каждого слагаемого внутри скобок на $-1$:

\[
-(a+b)=-a-b,
\qquad
-(a-b)=-a+b.
\]

\textbf{Примеры:}

\[
-(x-5)=-x+5,
\]

\[
-2(x-5)=-2x+10,
\]

\[
2(-x+5)=-2x+10,
\]

\[
-2(-x-5)=2x+10.
\]

\textbf{Контроль знаков:}

\begin{itemize}
  \item внешний множитель умножается на каждое
  слагаемое внутри скобок;

  \item отрицательный множитель меняет знак
  каждого произведения;

  \item слагаемое вне скобок сохраняет свой знак.
\end{itemize}

% ============================================================

\section*{3. Упрощение выражения}
\topicline

Для выражений этого типа последовательность действий такова:

\begin{enumerate}
  \item раскрыть все скобки;
  \item найти группы подобных слагаемых;
  \item привести подобные;
  \item проверить знаки и арифметику.
\end{enumerate}

Рассмотрим выражение

\[
-3(-1+5x)-5(10x-3).
\]

Раскрываем скобки:

\[
3-15x-50x+15.
\]

Приводим подобные:

\[
18-65x.
\]

В длинной строке полезно визуально группировать
одинаковые буквенные части:
отдельно слагаемые с $x$,
отдельно с $y$,
отдельно числа.

% ============================================================

\section*{4. Сначала упростить, затем подставить}
\topicline

Рассмотрим выражение

\[
5(2x-3y)-2(y-x)+3(2x-5y).
\]

После раскрытия скобок:

\[
10x-15y-2y+2x+6x-15y.
\]

После приведения подобных:

\[
18x-32y.
\]

При

\[
x=-1,
\qquad
y=-2
\]

этап подстановки записываем отдельной строкой:

\[
18\cdot(-1)-32\cdot(-2)
=
-18+64
=
46.
\]

Для всех выражений этого чек-листа

\[
\term{ОДЗ:}
\qquad
x,y,a,b,\ldots\in\mathbb{R},
\]

поскольку деления на выражения с переменной,
корней, логарифмов и других операций,
ограничивающих область допустимых значений,
здесь нет.

% ============================================================

\section*{5. Типичные ошибки и самопроверка}
\topicline

\begin{tabularx}
{\linewidth}
{@{}>{\raggedright\arraybackslash}p{0.31\linewidth}X@{}}

\toprule

\textbf{Ошибка}
&
\textbf{Как проверить}
\\

\midrule

Потерян знак при раскрытии скобок
&
Снова умножьте внешний множитель на каждое
слагаемое внутри скобок.
\\

Сложены неподобные слагаемые
&
Сравните буквенные части и степени переменных.
\\

Пропущено одно слагаемое
&
После раскрытия пересчитайте члены,
полученные из каждой скобки.
\\

Подстановка выполнена в уме
&
Запишите отдельной строкой выражение
с подставленными числами.
\\

Ошибка в $(-)\cdot(-)$
&
Перед вычислением отдельно определите
знак произведения.
\\

\bottomrule

\end{tabularx}

\vspace{0.7em}

\textbf{Мини-проверка перед ответом:}

\begin{enumerate}
  \item Все ли скобки раскрыты?

  \item Каждый ли внешний множитель применён
  ко всем членам скобки?

  \item Все ли подобные слагаемые собраны?

  \item Проверены ли знаки после раскрытия
  и подстановки?
\end{enumerate}

% ============================================================
% БЛОК-СХЕМА
% Отдельная страница.
%
% ВАЖНО:
% здесь намеренно нет глобальных TikZ-стилей вида
% flowproc/.style, flowcond/.style и т. п.
% Все параметры заданы непосредственно в узлах,
% поэтому подобные строки не могут попасть в вывод документа.
% ============================================================

\clearpage

\thispagestyle{fancy}

\begin{center}

{\Large\bfseries\color{darkaccent}
Алгоритм: упростить выражение и найти его значение
\par}

\vspace{1.2cm}

\begin{tikzpicture}[
  node distance=8mm
]

% ------------------------------------------------------------
% Начало
% ------------------------------------------------------------

\node[
  ellipse,
  draw=darkaccent,
  line width=0.9pt,
  minimum width=38mm,
  minimum height=10mm,
  align=center,
  inner sep=4pt
] (start)
{Начало};

% ------------------------------------------------------------
% Чтение выражения
% ------------------------------------------------------------

\node[
  below=of start,
  rounded corners=2mm,
  rectangle,
  draw=accent,
  line width=0.9pt,
  text width=66mm,
  minimum height=12mm,
  align=center,
  inner sep=5pt
] (read)
{Прочитайте выражение\\
и отметьте скобки, знаки, переменные};

% ------------------------------------------------------------
% Раскрытие скобок
% ------------------------------------------------------------

\node[
  below=of read,
  rounded corners=2mm,
  rectangle,
  draw=accent,
  line width=0.9pt,
  text width=66mm,
  minimum height=12mm,
  align=center,
  inner sep=5pt
] (open)
{Раскройте все скобки:\\
умножьте внешний множитель\\
на каждый член};

% ------------------------------------------------------------
% Подобные слагаемые
% ------------------------------------------------------------

\node[
  below=of open,
  rounded corners=2mm,
  rectangle,
  draw=accent,
  line width=0.9pt,
  text width=66mm,
  minimum height=12mm,
  align=center,
  inner sep=5pt
] (collect)
{Найдите подобные слагаемые\\
и приведите их};

% ------------------------------------------------------------
% Проверка необходимости подстановки
% ------------------------------------------------------------

\node[
  below=10mm of collect,
  diamond,
  aspect=2.25,
  draw=darkaccent,
  line width=0.9pt,
  text width=37mm,
  align=center,
  inner sep=2pt
] (needvalue)
{Нужно найти\\
числовое значение?};

% ------------------------------------------------------------
% Подстановка
% ------------------------------------------------------------

\node[
  right=24mm of needvalue,
  rounded corners=2mm,
  rectangle,
  draw=accent,
  line width=0.9pt,
  text width=56mm,
  minimum height=12mm,
  align=center,
  inner sep=5pt
] (subst)
{Подставьте значения\\
переменных\\
отдельной строкой и вычислите};

% ------------------------------------------------------------
% Финальная проверка
% ------------------------------------------------------------

\node[
  below=13mm of needvalue,
  rounded corners=2mm,
  rectangle,
  draw=accent,
  line width=0.9pt,
  text width=66mm,
  minimum height=12mm,
  align=center,
  inner sep=5pt
] (check)
{Проверьте раскрытие скобок,\\
коэффициенты, плюсы и минусы};

% ------------------------------------------------------------
% Ответ
% ------------------------------------------------------------

\node[
  below=of check,
  ellipse,
  draw=darkaccent,
  line width=0.9pt,
  minimum width=38mm,
  minimum height=10mm,
  align=center,
  inner sep=4pt
] (finish)
{Ответ};

% ------------------------------------------------------------
% Стрелки
% ------------------------------------------------------------

\draw[
  -{Stealth[length=3mm,width=2mm]},
  line width=0.9pt,
  color=textgray
]
(start) -- (read);

\draw[
  -{Stealth[length=3mm,width=2mm]},
  line width=0.9pt,
  color=textgray
]
(read) -- (open);

\draw[
  -{Stealth[length=3mm,width=2mm]},
  line width=0.9pt,
  color=textgray
]
(open) -- (collect);

\draw[
  -{Stealth[length=3mm,width=2mm]},
  line width=0.9pt,
  color=textgray
]
(collect) -- (needvalue);

% Ветка «да»

\draw[
  -{Stealth[length=3mm,width=2mm]},
  line width=0.9pt,
  color=textgray
]
(needvalue.east)
--
node[
  above,
  fill=white,
  inner sep=1pt,
  font=\small
]
{да}
(subst.west);

% Ветка «нет»

\draw[
  -{Stealth[length=3mm,width=2mm]},
  line width=0.9pt,
  color=textgray
]
(needvalue.south)
--
node[
  right,
  fill=white,
  inner sep=1pt,
  font=\small
]
{нет}
(check.north);

% Возврат от подстановки к проверке

\draw[
  -{Stealth[length=3mm,width=2mm]},
  line width=0.9pt,
  color=textgray
]
(subst.south)
|-
(check.east);

% Финальный переход

\draw[
  -{Stealth[length=3mm,width=2mm]},
  line width=0.9pt,
  color=textgray
]
(check) -- (finish);

\end{tikzpicture}

\end{center}

\clearpage

% ============================================================
% ТРЕНИРОВОЧНЫЙ БЛОК
% ============================================================

\section*{Тренировочный блок · Путь Б}
\topicline

Задачи расположены по нарастающей сложности.

Решения оформляйте в два этапа:

\[
\text{раскрытие скобок}
\;\longrightarrow\;
\text{приведение подобных}.
\]

В задачах на значение выражения этап подстановки
записывайте отдельной строкой.

% ============================================================
% 10.09.2026
% ============================================================

\subsection*{10.09.2026}

\begin{enumerate}[label=\textbf{\arabic*.}]

  \item Приведите подобные слагаемые:

  \[
  8a-11a+5b-2b.
  \]

  \item Раскройте скобки и упростите:

  \[
  6(x-4)-2(x+5).
  \]

  \item Раскройте скобки и упростите:

  \[
  -(3x-7)+2(2x+1).
  \]

  \item Упростите выражение:

  \[
  4(2x-3y)-3(y-x).
  \]

  \item Сначала упростите выражение,
  затем найдите его значение при

  \[
  x=-2,
  \qquad
  y=1:
  \]

  \[
  3(2x+y)-2(x-4y)+5(y-x).
  \]

\end{enumerate}

% ============================================================
% 11.09.2026
% ============================================================

\subsection*{11.09.2026}

\begin{enumerate}[label=\textbf{\arabic*.},resume]

  \item Приведите подобные слагаемые:

  \[
  -5m+2m+7n-9n.
  \]

  \item Раскройте скобки и упростите:

  \[
  -3(2x-5)+4(x-1).
  \]

  \item Упростите выражение:

  \[
  2(3x-y)-4(x+2y)+5(y-x).
  \]

  \item Упростите выражение:

  \[
  -(2x-3y)-2(4y-x)+3(x+y).
  \]

  \item Сначала упростите выражение,
  затем найдите его значение при

  \[
  x=-1,
  \qquad
  y=-2:
  \]

  \[
  4(2x-y)-3(x+2y)-2(3x-y).
  \]

\end{enumerate}

% ============================================================
% ОТВЕТЫ
% ============================================================

\clearpage

\section*{Ответы к тренировочному блоку}
\topicline

\begin{center}

\begin{tabularx}{\linewidth}{@{}c c X@{}}

\toprule

\textbf{Дата}
&
\textbf{№}
&
\textbf{Ответ}
\\

\midrule

10.09.2026
&
1
&
$-3a+3b$
\\

10.09.2026
&
2
&
$4x-34$
\\

10.09.2026
&
3
&
$x+9$
\\

10.09.2026
&
4
&
$11x-15y$
\\

10.09.2026
&
5
&
после упрощения:
$-x+16y$;
значение:
$18$
\\

\midrule

11.09.2026
&
6
&
$-3m-2n$
\\

11.09.2026
&
7
&
$-2x+11$
\\

11.09.2026
&
8
&
$-3x-5y$
\\

11.09.2026
&
9
&
$3x-2y$
\\

11.09.2026
&
10
&
после упрощения:
$-x-8y$;
значение:
$17$
\\

\bottomrule

\end{tabularx}

\end{center}

\end{document}